\section{Introduction}\label{sec:intro}

\subsection{Pseudorandomness from primes}
Fix a modulus $q\ge 1$ and write $\Gq=(\Z/q\Z)^\times$. A probability measure $\mu$ on positive
integers coprime to $q$ pushes forward to a probability measure on $\Gq$. We measure how
pseudorandom it is against the class of \emph{multiplicative characters} of $\Gq$, which are the
nontrivial Dirichlet characters $\chi$ modulo $q$:
\[
  \bias(\mu)=\max_{\chi\ne\chi_0 \bmod q}\Bigl\lvert\sum_n \mu(n)\chi(n)\Bigr\rvert .
\]
This is statistical pseudorandomness for a fixed test class. It is not cryptographic
indistinguishability. It is also a spectral quantity: the characters of $\Gq$ are eigenfunctions of the
averaging operator $Tf(g)=\sum_n\mu(n)f(gn)$ of the weighted Cayley graph $\mathrm{Cay}(\Gq,\mu)$, with
eigenvalues $\sum_n\mu(n)\chi(n)$ (Lemma~\ref{lem:eigen}). Small bias therefore means a uniform
spectral gap. Through the usual consequences of expansion (mixing of random walks, hitting and
Chernoff bounds for expander walks, coset hitting), it yields samplers on $\Gq$ that require
neither the factorization of $q$ nor knowledge of the group structure of $\Gq$.

Natural measures of this kind are supported on primes. Under the generalized Riemann hypothesis for
Dirichlet $L$-functions, the primes $p\le x$ form an expanding generating set of $\Gq$ once $x$ is a
little larger than $(\log q)^2$. This is the case $K=\mathbb{Q}$ of Jao, Miller and Venkatesan
\cite{jmv2009}, whose nontrivial eigenvalues are $O(x^{1/2}\log(xq))$ with absolute but unspecified
constants. Unconditionally, no polylogarithmic-support construction of this kind is known for $\Gq$.
One would imply a polylogarithmic bound for the least character non-residue, which is far beyond
current unconditional methods.

\subsection{The new input}
The Riemann hypothesis and GRH for all Dirichlet $L$-functions have recently been proved in Lean~4
\cite{alpoge2026,grhlean}. The proof has been independently built, replayed by several kernels and
compared against its statement files \cite{sutherlandgrh,sutherlandshortgrh,rubinsteinpage}. The Lean
development is the proof: the accompanying paper reduces the main inequality to stated hypotheses
rather than proving it \cite{rubinsteinpage}. The formal statement we use is
\[
  L(s,\chi)\ne0\quad\text{for all $q\ge1$, all $\chi$ mod $q$, and all } \operatorname{Re}s>\tfrac12,
\]
except at the pole $s=1$ of the principal character. Results that were conditional on GRH for
Dirichlet $L$-functions are thus unconditional, provided one accepts the Lean proof. This covers
neither Hecke $L$-functions of general number fields nor Artin $L$-functions, so the class-group and
isogeny-graph expanders of \cite{jmv2009} remain conditional.

\subsection{Results}
We prove explicit, fully formalized spectral bounds for prime Cayley graphs on $\Gq$.
\begin{enumerate}
\item \textbf{A uniform prime character-sum bound} (Theorem~\ref{thm:A}). For all $q\ge1$, all
  $\chi\ne\chi_0$ and all $N\ge3$,
  \[
    \Bigl\lvert\sum_{0<n<N}\chi(n)\Lambda(n)\bigl(1-\tfrac nN\bigr)\Bigr\rvert
    \le \frac{48e}{\pi}(\log N)^2\sqrt N\,\bigl(\pi\log(2+8q)+J\bigr),
  \]
  where $J=\int_\R\log(4+|t|)(1+t^2)^{-1}\,dt\le6$. The proof uses only nonvanishing of $L(s,\chi)$
  right of the critical line, Borel--Carath\'eodory and one contour shift. It needs no zero counting
  and no explicit formula.
\item \textbf{A sharp explicit bias bound} (Theorems~\ref{thm:sharp} and~\ref{thm:closed}). For a
  $\cosh$-based test function $f_a$ with kernel $k_a=f_a*\tilde f_a$, and weights
  $w(n)=\Lambda(n)k_a(\log n-y)n^{-1/2}$ on units $n\le e^{y+2a}$, we bound the bias by an explicit
  closed-form ratio $\mathrm{num}_B/\mathrm{den}_B$ for all $q$, $a\ge1$ and $y>2a$. The leading term
  is $(a+\sinh a)\log q/(\sqrt{x}\,(\sinh a'+a')^2)$ with $x=e^y$ and $a'=a-1/64$. The proof runs the
  Weil explicit formula for $L(s,\chi)$ and for $\zeta(s)$ through a two-sided positivity argument.
\item \textbf{A machine-checked certificate} (Corollary~\ref{cor:cert}). \emph{For every
  $1\le q\le2^{64}$ and every nontrivial $\chi$ modulo $q$, the measure above with $a=1$ and
  $y=2\log150$ has bias at most $1/2$.} It is supported on prime powers $n\le166253$ coprime to $q$.
  The same guarantee from Theorem~\ref{thm:A} would need $N\approx10^{16}$.
\item \textbf{Optimality of the kernel within the method} (Proposition~\ref{prop:opt}). Among all
  kernels $k=f*\tilde f$ with $\operatorname{supp}f\subseteq[-a,a]$, the leading coefficient
  $k(0)/\int k(u)e^{u/2}du$ produced by the positivity method is at least $1/(a+\sinh a)$, with
  equality exactly for $f\propto\cosh(t/2)\mathbf 1_{[-a,a]}$. Optimizing over $a$ gives
  $\bias\le(c^*+o(1))\log q/\sqrt X$ for support up to $X$, with $c^*=\min_a e^a/(a+\sinh a)=2/a_*^2\approx1.2236$
  attained at the root $a_*\approx1.2785$ of $a=1+e^{-a}$. No kernel gives a smaller leading constant
  within the positivity method.
\end{enumerate}

\subsection{What is new, and what is not}
The shape of Theorem~\ref{thm:A}, and of the leading term in Theorem~\ref{thm:sharp}, is classical
under GRH: compare Bach \cite{bach1990} and Lamzouri--Li--Soundararajan \cite{lls2015} for one-sided
statements, and \cite{jmv2009} for the spectral one. In particular, Lemma~6.1 of \cite{lls2015} already gives the leading term of
Theorem~\ref{thm:sharp} asymptotically (with a $1+o(1)$ factor) for each character, for a general class of kernels. Explicit GRH bounds for
primes in progressions, uniform in $q$, are in \cite{ehp2022}. We claim novelty for the following, with the caveat that a
literature search cannot establish uniqueness:
(i) bias (equivalently, second-eigenvalue) bounds for prime Cayley graphs on $\Gq$ with \emph{fully explicit}
constants and no $o(1)$ terms, uniform in $q$ and in $\chi$ (primitive or not), sharp enough to give small concrete certificates
such as Corollary~\ref{cor:cert};
(ii) the exact optimum of the leading constant within the positivity method (Proposition~\ref{prop:opt});
(iii) the end-to-end formal verification of (i) and of a concrete certificate, with GRH entering as a
proved theorem rather than an assumption.
The prior results closest to ours are discussed in Section~\ref{sec:related}. In particular, Bach's
constant~$2$ for the least non-residue is already formalized \cite{grhmiller}, and so is the Weil
explicit formula \cite{zeta23}. Our two-sided bounds use the same positivity inequality as
\cite{grhmiller}, in a different (spectral) way.

\subsection{Scope}
We do not improve the exponent $2$ in the $(\log q)^2$ support barrier (Section~\ref{sec:limits}):
doing so is at least as hard as improving the GRH bound for the least prime non-residue. We make no
cryptographic claims. Empirical statistics in Section~\ref{sec:experiments} are evidence about typical
behaviour; they are not proofs.
